\section{总结与延伸}

\subsection{作者的核心洞察}

\begin{importantbox}{Harness 设计的演化论}
\textit{``The space of interesting harness combinations doesn't shrink as models improve. Instead, it moves, and the interesting work for AI engineers is to keep finding the next novel combination.''}

有趣的 harness 组合空间不会随着模型改进而缩小——它会\textbf{移动}。AI 工程师的核心工作是持续寻找下一个有价值的组合。
\end{importantbox}

作者提炼的三个实践建议：

\begin{enumerate}
    \item \textbf{实验驱动}：用目标模型在真实问题上实验，阅读 trace，调优性能
    \item \textbf{按需分解}：对于更复杂的任务，将任务分解并对每个方面应用专门的 agent 可以带来额外收益
    \item \textbf{模型升级时重审 harness}：剥离不再承载性能的组件，添加新组件以达成之前不可能的能力
\end{enumerate}

\subsection{结构化提炼}

本文的贡献可以从三个层次理解：

\textbf{方法论层面}：建立了``假设-验证-迭代''的 harness 开发框架。每个 harness 组件都是一个可测试的假设（``模型不能独立完成 X''），当模型进步时假设需要重新验证。

\textbf{架构层面}：从 GAN 到 multi-agent 流水线的迁移展示了一个通用模式——将人类软件工程实践（代码审查、QA、sprint planning）编码为 agent 角色，而非试图让单个 agent 同时承担所有角色。

\textbf{工程实践层面}：
\begin{itemize}
    \item 评分标准本身就是有效的 prompt 工程——即使没有 evaluator 反馈，仅将评分标准写入 prompt 就能提升 baseline
    \item 文件作为 agent 间通信媒介是一种简单而有效的设计
    \item Planner 应保持高层次以避免错误级联，sprint contract 负责弥合细节间隙
    \item QA agent 的校准是一个需要多轮人工审查的调优过程，而非一次性配置
\end{itemize}

\subsection{拓展阅读}

\begin{itemize}
    \item Anthropic: \textit{Building Effective Agents} — harness 设计的基础原则
    \item Anthropic: \textit{Context Engineering} — 上下文窗口管理策略
    \item Claude Agent SDK 文档 — 构建 multi-agent 系统的官方框架
    \item Playwright MCP — 让 agent 与真实浏览器交互的工具
\end{itemize}

%% --- Fin del contenido --- %%

\end{document}
